% =============================================================================
% methods_formula_patches.tex
%
% Patch bundle for `methods_draft.tex` — converts the Al(Mg)-only methods to a
% generic A-host / B-solute framework that covers BOTH Al(Mg) and Pt(Au), and
% fills in the formula gaps a careful examiner will ask about (Δ E definition,
% X_GB, kJ/mol conversion, ceiling, Voronoi, anneal time constants, skew-normal,
% KS, a-CNA, block bootstrap, fit-parameter bootstrap CI).
%
% Each PATCH below is a self-contained LaTeX block. Paste-locations are stated
% at the top of each block. PATCH 0 also lists the find / replace edits needed
% in the existing prose to remove residual Al(Mg)-specific phrasing.
%
% Generated 2026-05-13 afternoon.
% =============================================================================


% =============================================================================
% PATCH 0a — Generic A/B convention + Systems-studied Table
% PASTE: at the very top of \section{Methodology}, REPLACING the first
%        paragraph's "global Mg fraction" → "global solute fraction" and
%        adding the convention sentence + table reference.
% =============================================================================

Our objective is to test where the independent-site Fermi--Dirac (FD)
framework of Wagih \emph{et al.}~\cite{Wagih2020} for grain-boundary (GB)
solute segregation breaks down in $(T, X_c)$ space, where $T$ is the
temperature and $X_c$ the global solute fraction. The methodology is
structured as a hypothesis-versus-ground-truth comparison on a single
shared atomistic basis (Fig.~\ref{fig:method}): the FD framework provides
the prediction $X_\mathrm{GB}^\mathrm{FD}(T, X_c)$, hybrid Monte Carlo /
molecular dynamics (HMC) sampling provides the ground truth
$X_\mathrm{GB}^\mathrm{HMC}(T, X_c)$, and the gap between them is the
testable quantity. Throughout this work we denote the host element $A$ and
the dilute solute $B$, and report results for two binary systems
summarised in Table~\ref{tab:systems}. All equations and protocols in
this section apply to both systems; values that differ between systems
($N_\mathrm{tot}$, $N_\mathrm{GB}$, $T_\mathrm{melt}^{(A)}$, etc.) are
tabulated rather than hard-coded in the text. The equations and the HMC
algorithm are stated in Apps.~\ref{app:eq}--\ref{app:alg}.


% =============================================================================
% PATCH 0b — Systems-studied Table
% PASTE: immediately after PATCH 0a's paragraph, before \subsection{Atomistic
%        base...}.
% =============================================================================

\begin{table}[t]
\centering
\caption{Binary systems studied. Both are FCC hosts with a dilute
substitutional solute, processed by the same generic pipeline
(Fig.~\ref{fig:method}); only the system-specific quantities below
differ. The bulk modulus column motivates the relative CG-iteration
cost reported in App.~\ref{app:dE_compare}. Note that $f_\mathrm{GB}$
is geometry- and size-dependent (smaller boxes have a higher
surface-to-volume ratio so a larger fraction of atoms sits in GB
environments); the higher $f_\mathrm{GB}$ for Pt(Au) reflects its
smaller box edge $L = 100$~\AA{} rather than an intrinsic material
property.}
\label{tab:systems}
\begin{tabular}{lcc}
\toprule
                                          & Al(Mg)            & Pt(Au)              \\
\midrule
host $A$                                  & Al                & Pt                  \\
solute $B$                                & Mg                & Au                  \\
EAM potential                             & Mendelev 2009~\cite{Mendelev2009}      & O'Brien 2017~\cite{OBrien2017} \\
\quad LAMMPS \texttt{pair\_style}         & \texttt{eam/fs}                        & \texttt{eam/alloy}             \\
lattice constant $a$ (\AA)                & 4.05              & 3.98                \\
box edge $L$ (\AA)                        & 200               & 100                 \\
total atoms $N_\mathrm{tot}$              & 475{,}715         & 62{,}096            \\
grain seeds $N_\mathrm{grains}$           & 16                & 8                   \\
GB atoms $N_\mathrm{GB}$ (post-anneal)    & 89{,}042          & 23{,}272            \\
GB fraction $f_\mathrm{GB}$               & 0.187             & 0.375               \\
host melting $T_\mathrm{melt}^{(A)}$ (K)  & 933               & 2041                \\
anneal hold $T_\mathrm{hold} = 0.4 T_\mathrm{melt}^{(A)}$ (K) & 373 & 816 \\
$\Delta E_i$ spectrum sample $n$          & 500               & 500                 \\
host bulk modulus $K$ (GPa)               & 76                & 280                 \\
\bottomrule
\end{tabular}
\end{table}


% =============================================================================
% PATCH 1 — ΔE simplification note
% PASTE: in app:eq, immediately AFTER equation eq:dE (within the same paragraph
%        or as a follow-on paragraph).
% =============================================================================

For a pure-$A$ polycrystal base the bracketed Al-substituted terms
satisfy $E^{\mathrm{bulk-}A} = E_i^{\mathrm{GB-}A} = E_\mathrm{pure}$
(both equal the total energy of the un-substituted polycrystal), and
Eq.~(\ref{eq:dE}) simplifies to the operational two-point form
\begin{equation}\label{eq:dE_op}
\Delta E_i \;=\; E_i^{\mathrm{GB-}B} \;-\; \big\langle E^{\mathrm{bulk-}B} \big\rangle
\qquad (\text{pure-}A\ \text{host}),
\end{equation}
in which $E_i^{\mathrm{GB-}B}$ is the total potential energy of the
polycrystal with a single $A \to B$ substitution at GB site $i$
followed by conjugate-gradient minimisation with tight tolerances
$\mathrm{etol} = \mathrm{ftol} = 10^{-25}$ (matching
Ref.~\cite{Wagih2020}'s \texttt{calculate\_E\_GB\_solute.in} protocol),
and $\langle E^{\mathrm{bulk-}B} \rangle$ is the same total energy
averaged over $n_\mathrm{ref} = 10$ well-separated bulk-$A$ reference
positions. Eq.~(\ref{eq:dE_op}) is what
\texttt{scripts/sample\_delta\_e.py} evaluates per GB site; the bulk
reference scatter ($\sigma_{E^{\mathrm{bulk-}B}} = 8.1$~meV for Al(Mg);
$5.3$~meV for Pt(Au)) is one to two orders of magnitude below typical
per-site signal magnitudes ($|\Delta E_i| \sim 100$--$500$~meV). We
sample $n = 500$ random GB sites per system (Table~\ref{tab:systems}),
which is sufficient to fit the skew-normal of
Eq.~(\ref{eq:skewnorm}) below to the precision required by the
canonical FD predictor.


% =============================================================================
% PATCH 2 — X_GB(t) definition (currently undefined despite being everywhere)
% PASTE: in app:eq, after the canonical FD equation eq:FD_canon (define
%        before the HMC equations reference it).
% =============================================================================

\paragraph*{HMC observable.}
The instantaneous GB solute fraction at production time $t$ is
\begin{equation}\label{eq:Xgb_def}
X_\mathrm{GB}(t) \;\equiv\; \frac{1}{N_\mathrm{GB}}
\sum_{i \in \mathrm{GB}}
\mathbb{1}\!\big\{\, \text{site $i$ is occupied by $B$ at time $t$} \,\big\} ,
\end{equation}
where $\mathbb{1}\{\cdot\}$ is the indicator function and the sum is
over the $N_\mathrm{GB}$ atoms tagged as GB by the a-CNA mask
(App.~\ref{app:atomistic}). The reported HMC observable is the
post-burn-in time average $\langle X_\mathrm{GB}\rangle_t$ together
with the stationary block-bootstrap 95\% confidence interval defined
in Sec.~\ref{sec:meth:stats}, Eq.~(\ref{eq:block_bootstrap}).


% =============================================================================
% PATCH 3 — Unit conversion (kJ/mol ↔ eV) and closed-box ceiling
% PASTE: in app:eq, at the top (these are conversion / bound formulas the
%        reader needs before the substantive equations).
% =============================================================================

\paragraph*{Energy units.}
LAMMPS evaluates total energies in eV (metal units). All $\Delta E_i$
histograms, fits, and figures throughout this work are reported in
kJ/mol via the standard conversion
\begin{equation}\label{eq:unit_kjmol}
\Delta E \;[\mathrm{kJ/mol}]
\;=\; \frac{e \cdot N_A}{1000} \, \Delta E \;[\mathrm{eV}]
\;=\; 96.485 \, \Delta E \;[\mathrm{eV}] ,
\end{equation}
where $e = 1.602 \times 10^{-19}$~C is the elementary charge and
$N_A = 6.022 \times 10^{23}\,\mathrm{mol}^{-1}$ is Avogadro's number
(the factor $10^{-3}$ converts J/mol to kJ/mol).

\paragraph*{Closed-box ceiling on $X_\mathrm{GB}$.}
A closed simulation cell with global solute fraction $X_c$ has a hard
upper bound on $X_\mathrm{GB}$ set by mass conservation alone — the
configuration with every solute atom piled at a GB site:
\begin{equation}\label{eq:ceiling}
X_\mathrm{GB}^{\mathrm{ceil}}(X_c) \;=\; X_c \, N_\mathrm{tot} / N_\mathrm{GB}
\;=\; X_c / f_\mathrm{GB} .
\end{equation}
Both $X_\mathrm{GB}^{\mathrm{HMC}}$ and $X_\mathrm{GB}^{\mathrm{FD,canon}}$
are bounded above by Eq.~(\ref{eq:ceiling}); the ceiling line is
plotted alongside HMC and FD curves in the panel-(d) headline figure
of the Results.


% =============================================================================
% PATCH 4 — Voronoi region (replace the inline "3D Voronoi tessellation..."
%           in app:atomistic with the explicit formula).
% =============================================================================

\paragraph*{Voronoi tessellation.}
Given $N_\mathrm{grains}$ seed positions
$\{\mathbf{s}_k\}_{k=1}^{N_\mathrm{grains}}$ drawn uniformly inside the
periodic $L^3$ box, the grain assigned to each candidate FCC lattice
site $\mathbf{r}$ is
\begin{equation}\label{eq:voronoi}
k(\mathbf{r}) \;=\; \arg\min_{1 \le k \le N_\mathrm{grains}}
\big\| \mathbf{r} - \mathbf{s}_k \big\|_\mathrm{PBC} ,
\end{equation}
where $\|\cdot\|_\mathrm{PBC}$ is the minimum-image Euclidean distance
under periodic boundary conditions, implemented via a $k$-d tree
nearest-neighbour lookup. Each grain is given a uniformly random
$\mathrm{SO}(3)$ orientation (Haar measure on rotation matrices) before
the seed lattice is rotated into place. Lattice sites whose nearest
inter-grain distance is below the close-pair cutoff
$r_\mathrm{cut} = a/(2\sqrt{2})$ (one half the nearest-neighbour
spacing) are pruned to prevent unphysically close atomic pairs at
as-built GBs; the resulting deficit is $\sim 2\%$ of the ideal FCC
atom count (Table~\ref{tab:systems}).


% =============================================================================
% PATCH 5 — Anneal thermostat / barostat time constants
% PASTE: in app:atomistic, replacing or supplementing the existing inline
%        "anneal schedule follows Wagih..." description.
% =============================================================================

\paragraph*{Anneal schedule (Wagih protocol, generic).}
After the close-pair-pruned Voronoi lattice is assembled (App.~%
\ref{app:atomistic} above), the box is annealed in five stages:
\begin{enumerate}\itemsep0pt\parsep0pt\setlength{\leftmargin}{0pt}
\item Conjugate-gradient minimisation to absorb residual close-pair strain from the Voronoi construction (\texttt{etol} / \texttt{ftol} $= 10^{-6} / 10^{-8}$~eV / (eV/\AA), max iterations $5000$).
\item NVT temperature ramp from 1~K to $T_\mathrm{hold} = 0.4 T_\mathrm{melt}^{(A)}$ over 5~ps using a Nos\'e--Hoover thermostat with damping $\tau_T = 100$~fs.
\item NPT hold at $T_\mathrm{hold}$ for 250~ps with isotropic Nos\'e--Hoover thermo-barostat ($\tau_T = 100$~fs, $\tau_P = 1000$~fs, target $P = 0$) to relax GB structure.
\item NPT cool from $T_\mathrm{hold}$ to 1~K at $3$~K/ps with the same barostat settings.
\item Final zero-pressure box relaxation by conjugate gradient (\texttt{fix box/relax} \texttt{iso 0.0}; tolerances $10^{-8}$ / $10^{-10}$).
\end{enumerate}
The annealed configuration is the fixed substrate for all subsequent
measurements. Adaptive common-neighbour analysis~\cite{Stukowski2012}
classifies each atom (App.~PATCH~8) and the resulting GB mask is held
fixed throughout.


% =============================================================================
% PATCH 6 — Skew-normal PDF + moment-to-parameter relations
% PASTE: in app:dE_compare, at the top of the appendix (before the n=500
%        justification paragraph that already exists).
% =============================================================================

\paragraph*{Skew-normal fit.}
The $\Delta E_i$ spectrum is fit by the three-parameter skew-normal
PDF~\cite{Azzalini1985} (the same form used by Wagih's
SI~\cite{Wagih2020}):
\begin{equation}\label{eq:skewnorm}
f(\Delta E; \mu, \sigma, \alpha)
\;=\; \frac{2}{\sigma} \, \varphi\!\left(\tfrac{\Delta E - \mu}{\sigma}\right)
\, \Phi\!\left(\alpha \, \tfrac{\Delta E - \mu}{\sigma}\right),
\end{equation}
where $\varphi$ and $\Phi$ are the standard-normal PDF and CDF.
$\mu$ is the location parameter, $\sigma > 0$ the scale parameter,
and $\alpha$ the shape (skewness) parameter; $\alpha = 0$ recovers
the symmetric Gaussian. Parameters are obtained by maximum likelihood
(\texttt{scipy.stats.skewnorm.fit}). The sample mean and standard
deviation are NOT $\mu$ and $\sigma$ when $\alpha \neq 0$:
\begin{equation}\label{eq:skewnorm_moments}
\langle \Delta E \rangle = \mu + \sigma \delta \sqrt{2/\pi},
\qquad
\mathrm{var}(\Delta E)  = \sigma^2 \!\left(1 - \tfrac{2 \delta^2}{\pi}\right),
\qquad
\delta = \alpha / \sqrt{1+\alpha^2}.
\end{equation}
Fit parameters are reported in kJ/mol per the conversion
(\ref{eq:unit_kjmol}) and matched against Wagih's SI tables for each
$(A, B)$ pair.


% =============================================================================
% PATCH 7 — Two-sample KS statistic and its small-sample interpretation
% PASTE: in app:dE_compare, after PATCH 6.
% =============================================================================

\paragraph*{Spectrum-level comparison: two-sample Kolmogorov--Smirnov.}
Let $\hat F_1, \hat F_2$ denote the empirical CDFs of two $\Delta E$
samples of sizes $n_1, n_2$. The two-sample Kolmogorov--Smirnov
statistic and asymptotic two-sided $p$-value are
\begin{equation}\label{eq:KS}
D \;=\; \sup_{\Delta E} \big|\, \hat F_1(\Delta E) - \hat F_2(\Delta E) \,\big|,
\qquad
p \;=\; 2 \sum_{k=1}^{\infty} (-1)^{k-1} \exp(-2 k^2 \lambda^2),
\quad
\lambda = D \sqrt{\tfrac{n_1 n_2}{n_1 + n_2}},
\end{equation}
implemented in \texttt{scipy.stats.ks\_2samp}. We report
$(D, p)$ for our $n = 500$ spectrum versus each Wagih Zenodo reference
($n_2 = 82{,}646$ for Al(Mg), $n_2 = 97{,}440$ for Pt(Au)). The KS
$p$-value alone is sample-size-asymmetric: at $n_1 \ll n_2$ a small
$p$ may reflect tail under-sampling on the small-$n$ side rather than
a real distribution mismatch (Pt(Au) yields $p = 6 \times 10^{-4}$ for
this reason). We therefore supplement the KS test with the
fit-parameter bootstrap CI of Eq.~(\ref{eq:fit_bootstrap}) as the
primary spectrum-validation metric.


% =============================================================================
% PATCH 8 — a-CNA brief description (replaces the existing one-liner in
%           app:atomistic about "Adaptive common-neighbour analysis with
%           default cutoffs").
% =============================================================================

\paragraph*{GB identification.}
The annealed configuration is classified atom-by-atom using adaptive
common-neighbour analysis (a-CNA)~\cite{Stukowski2012}, which assigns
each atom one of $\{\text{FCC}, \text{BCC}, \text{HCP}, \text{ICO},
\text{other}\}$ based on the topology of its $12$ nearest neighbours.
The cutoff radius for the neighbour set is determined adaptively per
atom (hence ``adaptive''), avoiding fixed-cutoff failures near
distorted regions. Atoms not classified as FCC (the bulk lattice
phase) are designated \emph{GB atoms}. The resulting binary mask
of $N_\mathrm{GB}$ sites (Table~\ref{tab:systems}) is generated once
on the annealed pure-$A$ polycrystal and held fixed across all
downstream measurements — the $\Delta E$ sampling, the FD predictor
[via $f_\mathrm{GB}$ in Eq.~(\ref{eq:ceiling}) and the summand of
Eq.~(\ref{eq:FD_canon})], and the HMC observable
(Eq.~\ref{eq:Xgb_def}).


% =============================================================================
% PATCH 9 — Stationary block bootstrap on ⟨X_GB⟩ trajectory mean
% PASTE: in sec:meth:stats, replacing the existing one-sentence description
%        ("stationary block bootstrap with a 5-frame block size and 2000
%        resamples").
% =============================================================================

\paragraph*{Trajectory-level CI on $\langle X_\mathrm{GB}\rangle_t$.}
Let $\{x_t\}_{t=1}^{n}$ be the post-burn-in $X_\mathrm{GB}(t)$ frames
sampled at 1~ps spacing (the first 20\% of frames are dropped as
burn-in; $n \approx 240$ for our default 300~ps production). The
stationary block bootstrap of Politis \& Romano~\cite{Politis1994}
generates $B = 2000$ replicate samples $\{x_t^{(b)}\}_{t=1}^{n}$ by
concatenating blocks of random geometric length with mean
$L_\mathrm{block} = 5$ frames, sampled with replacement from
$\{x_t\}$. Each replicate yields a bootstrap mean $\bar x^{(b)}$, and
the 95\% CI is the empirical $[2.5, 97.5]$ percentile range of
$\{\bar x^{(b)}\}_{b=1}^{B}$:
\begin{equation}\label{eq:block_bootstrap}
\mathrm{CI}_{95\%}\!\left(\langle X_\mathrm{GB}\rangle_t\right)
\;=\; \Big[\, q_{2.5}\!\left(\{\bar x^{(b)}\}\right), \;
              q_{97.5}\!\left(\{\bar x^{(b)}\}\right) \,\Big].
\end{equation}
The geometric block length is chosen large enough to span the
autocorrelation time of the $X_\mathrm{GB}(t)$ series under the swap
schedule of Algorithm~\ref{alg:hmc} (the dominant autocorrelation is
shorter than 5~ps in all converged runs, so $L_\mathrm{block} = 5$
frames over-covers).


% =============================================================================
% PATCH 10 — Fit-parameter bootstrap CI (the Pt(Au)-driven validation method,
%             complementary to KS when n_1 ≪ n_2).
% PASTE: in app:dE_compare, after PATCH 7.
% =============================================================================

\paragraph*{Fit-parameter bootstrap CI.}
When the Wagih reference contains $n_2 \gg n_1$ sites
[Eq.~(\ref{eq:KS})], the KS $p$-value becomes acutely sensitive to
small tail offsets that any random $n_1$ sub-sample of the larger pool
also exhibits. As a complementary spectrum-validation test we
therefore build the bootstrap distribution of the fit parameters
$\{\mu, \sigma, \alpha\}$ that one would obtain from random
$n_1$-sized sub-samples of the Wagih pool. Concretely:
\begin{enumerate}\itemsep0pt\parsep0pt
\item Draw $B = 10{,}000$ resamples
      $\{\Delta E_j^{(b)}\}_{j=1}^{n_1}$ with replacement from the
      Wagih reference pool of size $n_2$.
\item Fit each resample to Eq.~(\ref{eq:skewnorm}), yielding
      $\big(\mu^{(b)}, \sigma^{(b)}, \alpha^{(b)}\big)$.
\item Form the 95\% percentile CI in each parameter:
\begin{equation}\label{eq:fit_bootstrap}
\mathrm{CI}_{95\%}(\theta) \;=\; \Big[\, q_{2.5}\!\left(\{\theta^{(b)}\}_{b=1}^{B}\right), \;
                                              q_{97.5}\!\left(\{\theta^{(b)}\}_{b=1}^{B}\right) \,\Big],
\quad \theta \in \{\mu, \sigma, \alpha\}.
\end{equation}
\end{enumerate}
Our $n_1 = 500$ measurement is declared statistically consistent with
the Wagih reference if our maximum-likelihood
$(\hat\mu, \hat\sigma, \hat\alpha)$ all fall inside their bootstrap
CI95. For Pt(Au) the three fit parameters pass this test (with
$|z| \le 1.6$ on each) while the raw sample mean and standard
deviation fall outside their bootstrap CI95 at $z \approx \pm 2.7$;
the shape statistics are preserved while the deepest tail outliers
are not, indicating that our $n_1 = 500$ sample misses
$\Delta E < -34$~kJ/mol GB sites that exist in the larger Wagih pool
($\min \Delta E_{\mathrm{Wagih}} \approx -46$~kJ/mol). Because the FD
predictor [Eq.~(\ref{eq:FD_canon})] depends only on the
parametric form of the spectrum and not on its extreme-tail outliers,
the bootstrap-fit-CI is the operationally relevant
spectrum-validation criterion for this work.


% =============================================================================
% GENERICIZATION CHECKLIST — find/replace edits needed in EXISTING methods text
% to remove residual Al(Mg)-only phrasing once the patches above are in place.
% =============================================================================
%
% Locate                                              Replace with
% ------------------------------------------          --------------------------------------------
% "global Mg fraction"                                "global solute fraction"
% "200 Å cubic Al polycrystal"                        "$L^3$ cubic $A$ polycrystal (Table~\ref{tab:systems})"
% "N_tot = 475,715 atoms"                             "$N_\mathrm{tot}$ atoms (Table~\ref{tab:systems})"
% "16 grain seeds"                                    "$N_\mathrm{grains}$ grain seeds"
% "the Mendelev Al--Mg embedded-atom-method
%  potential~\cite{Mendelev2009}"                     "the alloy-specific EAM potential (Table~\ref{tab:systems})"
% "N_GB = 89,042"                                     "$N_\mathrm{GB}$ (Table~\ref{tab:systems})"
% "f_GB = 0.187"                                      "$f_\mathrm{GB}$ (Table~\ref{tab:systems})"
% "single Mg from a bulk substitutional site to GB"   "single $B$ atom from a bulk site to GB"
% "Mg ↔ Al type-swap moves"                           "$A \leftrightarrow B$ type-swap moves"
% "preserves the closed-box N_Mg"                     "preserves the closed-box $N_B$"
% "place ⌊X_GB^FD · N_GB⌋ Mg atoms at random GB"     "place $\lfloor X_\mathrm{GB}^\mathrm{FD} N_\mathrm{GB}\rfloor$ $B$ atoms at random GB"
% "remaining N_Mg = X_c · N_tot Mg atoms"             "remaining $N_B = X_c N_\mathrm{tot}$ $B$ atoms"
% "local Mg neighbour count n^Mg_local"               "local $B$-neighbour count $n^{B}_\mathrm{local}$"
% "Mg atoms within r ≤ 5 Å"                           "$B$ atoms within $r \leq 5$~\AA"
% "T_melt = 933 K" (anywhere)                          "$T_\mathrm{melt}^{(A)}$ (Table~\ref{tab:systems})"
% "0.4 T_melt = 373 K"                                "$T_\mathrm{hold} = 0.4 T_\mathrm{melt}^{(A)}$"
% "lattice constant is a = 4.05 Å"                    "$a$ (Table~\ref{tab:systems})"
% =============================================================================
